\begin{center}
    \section*{\fontsize{16}{16}\selectfont Chapter 1: Introduction}
\end{center}
\vspace{-2mm}

\section{Introduction}
Healthcare delivery in developing nations faces significant structural barriers, including fragmented workflows across symptom triage, provider discovery, consultation, and post-care record management. Consequently, patients frequently delay or forgo care. While standalone tools for triage, video calling, and payments exist, unified digital ecosystems remain absent. 

\noindent This work introduces \textbf{EHealth}, an integrated AI-assisted telemedicine platform offering an end-to-end care continuum: intelligent triage, provider scheduling, WebRTC video consultation, automated clinical documentation, and blockchain-verified prescription security.

\subsection{Problem Statement}
Digital healthcare systems suffer from three primary structural challenges:
\begin{enumerate}
    \setlength{\itemsep}{1pt}
    \setlength{\parskip}{0pt}
    \item \textbf{Triage Disconnect:} Patients lack objective guidance to evaluate symptom severity, leading to emergency resource misallocation or delayed care.
    \item \textbf{Documentation Overhead:} Manual clinical note generation places heavy administrative burdens on doctors, introducing documentation errors and variance.
    \item \textbf{Prescription Tampering:} Digital prescriptions lack independent, out-of-band verification frameworks, leaving files vulnerable to post-issuance modification.
\end{enumerate}

\subsection{Problem Background}
Centralized urban healthcare models leave semi-urban and rural populations under-resourced. Although post-COVID-19 telemedicine adoption grew, existing services function mostly as basic video wrappers without decision-support intelligence. For example, traditional telehealth systems depend heavily on manual documentation by doctors, which cuts into actual patient interaction time and increases fatigue—an issue that recent automated speech recognition systems aim to resolve \cite{asr_healthcare_2026, asr_llm_benchmarking_2025}. 

\noindent While some platforms have started using large language models for initial symptom checking and triage \cite{triage_llm_2024, inclusive_healthcare_2025}, these tools are usually completely separate from the actual video consultation and post-care record management. Patients end up juggling multiple disjointed apps just to get treated. On top of that, standard digital prescriptions remain vulnerable to unauthorized changes after they are issued, a security gap that typical centralized databases have struggled to close without the integration of smart contracts \cite{medical_blockchain_ehr_2022, blockchain_healthcare_2020}.

\noindent Our work, EHealth, takes a different approach by building these features directly into a single, unified workflow. Instead of fragmented steps, we bring everything together. Large Language Models (LLMs) handle the initial urgency classification before the patient even sees a doctor. During the call, lightweight speech recognition engines (like Whisper) run locally to transcribe the conversation, removing the heavy documentation burden that older models suffer from. To secure the medical records, high-throughput Proof-of-Stake (PoS) blockchains like Polygon make on-chain document hashing cheap and fast, directly solving the prescription tampering issues seen in conventional systems. Combining these technologies via WebRTC creates a robust, end-to-end clinical experience that stands in sharp contrast to the disconnected tools currently available.

\subsection{Research Objectives}
The primary objective is to design, deploy, and evaluate an integrated AI-driven telemedicine platform. Specific goals include:
\begin{itemize}
    \setlength{\itemsep}{1pt}
    \setlength{\parskip}{0pt}
    \item \textbf{Triage Module:} Build an LLM engine to classify unstructured text/voice inputs into urgency levels (\texttt{HIGH}, \texttt{MEDIUM}, \texttt{LOW}) using constrained JSON schemas.
    \item \textbf{Documentation Pipeline:} Construct a local, containerized pipeline to transcribe consultation audio and synthesize structured clinical notes off-line.
    \item \textbf{Prescription Integrity:} Deploy a Polygon PoS smart contract interface to anchor SHA-256 prescription hashes for dual-state validation.
    \item \textbf{Full Integration:} Unify triage, scheduling, WebRTC streaming, AI notes, and blockchain verification into a synchronized NestJS/Flutter stack.
    \item \textbf{Empirical Evaluation:} Quantify triage accuracy, ASR Word Error Rate (WER), clinical extraction F1-scores, and blockchain gas costs.
\end{itemize}

\subsection{Motivation}
The primary drive behind EHealth stems from the glaring disparities in modern healthcare delivery, especially in under-resourced regions where patients often delay seeking help due to fragmented and disjointed systems. While cutting-edge technologies like artificial intelligence and blockchain hold massive potential, they are frequently isolated in theoretical research or standalone applications. Our motivation was to bridge this gap by taking these advanced concepts and weaving them into a practical, everyday tool that directly improves how patients access care and how doctors manage their workload.

\subsection{Significance of Research}
Systemically, this project proves that a complex, 5-stage pipeline—spanning from AI triage to WebRTC consultations and local transcription—can operate seamlessly under a unified backend without suffering from compounding end-to-end latency. Economically and practically, it demonstrates that using public Proof-of-Stake blockchains to secure clinical artifacts is not just a novelty, but a highly feasible and scalable solution for medical record integrity. By addressing these critical engineering bottlenecks, our work lays a solid foundation for deploying integrated, next-generation telemedicine platforms in the real world.

\subsection{Research Contributions}
The core contributions of this project include:
\begin{itemize}
    \setlength{\itemsep}{1pt}
    \setlength{\parskip}{0pt}
    \item \textbf{Constrained Triage Engine:} Implemented a Redis-cached endpoint using Baichuan-M2-32B returning structured JSON triage schemas from noisy inputs.
    \item \textbf{Privacy-Preserving Documentation:} Built an offline pipeline (FFmpeg $\rightarrow$ \textit{Xenova/whisper-tiny} $\rightarrow$ LLM restructuring $\rightarrow$ PDF engine $\rightarrow$ Google Drive API) running entirely within local containers.
    \item \textbf{Dual-State Integrity Module:} Developed a Polygon PoS smart contract storing \texttt{keccak256} digests of SHA-256 hashes, exposing an endpoint for concurrent database and on-chain verification.
    \item \textbf{Full-Stack Infrastructure:} Integrated mediasoup SFU WebRTC routing, Socket.IO signaling, SSLCommerz payment processing, and JWT security into a unified API.
    \item \textbf{Quantitative Evaluation:} Benchmarked system modules, achieving 87.5\% triage accuracy (200 cases), 93.0\% extraction F1-score (50 dialogues), and a median write cost of 0.00607\,POL per record update.
\end{itemize}

\subsection{Complex Engineering Analysis}

\subsubsection{Complex Engineering Problems}
\begin{itemize}
    \setlength{\itemsep}{1pt}
    \setlength{\parskip}{0pt}
    \item \textbf{Unstructured Input Robustness:} Handling colloquial and fragmented patient text required explicit prompt schemas, temperature control, and fallback parsing branches.
    \item \textbf{Acoustic Audio Normalization:} Processing noisy WebRTC audio necessitated a deterministic FFmpeg preprocessing phase (16\,kHz mono WAV) prior to local ONNX speech inference.
    \item \textbf{Cryptographic State Alignment:} Ensuring exact hash parity between client verification uploads and on-chain state required strict byte-level normalization across MySQL and Solidity interfaces.
\end{itemize}

\subsubsection{Complex Engineering Activities}
\begin{itemize}
    \setlength{\itemsep}{1pt}
    \setlength{\parskip}{0pt}
    \item \textbf{Relational Database Normalization:} Designed a unified schema linking users, doctors, service fees, payment entities, and multi-version prescription transaction hashes.
    \item \textbf{SFU Media Containerization:} Configured a containerized mediasoup SFU balancing WebRTC transport negotiation with client-side audio buffering for downstream queues.
    \item \textbf{Cross-Runtime C-Library Fixes:} Resolved container build conflicts between Alpine Linux C-libraries (\texttt{musl}) and ONNX runtime requirements (\texttt{glibc}).
\end{itemize}

\subsection{Project Organization}
Chapter 2 reviews related literature. Chapter 3 presents the overall architecture. Chapter 4 details module methodologies. Chapter 5 covers system implementation. Chapter 6 presents quantitative evaluation results, followed by conclusions and future work in Chapter 7.

\vspace{2mm}
\noindent\textbf{Summary:} This chapter established the scope, architecture, and core engineering contributions of EHealth. The subsequent chapters detail the design, execution, and empirical evaluation of these subsystems.